% from ega3-1-fr.tex
\oldpage[III]{88}
est commutatif, on en conclut que les diagrammes
\begin{equation}
\label{III.1.2.5.1-fr}
\tag{1.2.5.1}
\xymatrix{
H^p(\mathfrak U,\mathcal F'')\ar[r]^{\partial}\ar[d]
&H^{p+1}(\mathfrak U,\mathcal F')\ar[d]\\
H^{p+1}((\mathbf f),M'')\ar[r]^{\partial}
&H^{p+2}((\mathbf f),M')
}
\end{equation}
sont commutatifs pour tout $p$ (G, I, 2.1.1).
\end{remark}

\subsection{Cohomologie d'un schéma affine}
\label{subsection:III.1.3-fr}

\begin{theorem}[1.3.1]
\label{III.1.3.1-fr}
Soit $X$ un schéma affine. Pour tout $\mathcal O_X$-Module quasi-cohérent
$\mathcal F$, on a $H^p(X,\mathcal F)=0$ pour tout $p>0$.
\end{theorem}

\begin{proof}
Soit $\mathfrak U$ un recouvrement fini de $X$ par des ouverts affines
$X_{f_i}=D(f_i)$ ($1\leq i\leq r$)~; on sait qu'alors l'idéal engendré dans
$A=\Gamma(X,\mathcal O_X)$ par les $f_i$ est égal à $A$. On conclut donc de
(\hyperref[III.1.2.4-fr]{1.2.4}) que l'on a
$H^p(\mathfrak U,\mathcal F)=0$ pour $p>0$. Comme il y a des recouvrements
finis de $X$ par des ouverts affines qui sont arbitrairement fins (I, 1.1.10),
la définition de la cohomologie de Čech (G, II, 5.8) montre que l'on a aussi
$\check H^p(X,\mathcal F)=0$ pour $p>0$. Mais ceci s'applique aussi à tout
préschéma $X_f$ pour $f\in A$ (I, 1.3.6), donc
$\check H^p(X_f,\mathcal F)=0$ pour $p>0$. Comme l'on a
$X_f\cap X_g=X_{fg}$, on en déduit que l'on a aussi
$H^p(X,\mathcal F)=0$ pour tout $p>0$, en vertu de (G, II, 5.9.2).
\end{proof}

\begin{corollary}[1.3.2]
\label{III.1.3.2-fr}
Soient $Y$ un préschéma, $f:X\to Y$ un morphisme affine (II, 1.6.1). Pour
tout $\mathcal O_X$-Module quasi-cohérent $\mathcal F$, on a
$R^qf_*(\mathcal F)=0$ pour $q>0$.
\end{corollary}

\begin{proof}
En effet, par définition $R^qf_*(\mathcal F)$ est le $\mathcal O_Y$-Module
associé au préfaisceau $U\mapsto H^q(f^{-1}(U),\mathcal F)$, $U$ parcourant
les ouverts de $Y$. Mais les ouverts affines forment une base de $Y$, et pour
un tel ouvert $U$, $f^{-1}(U)$ est affine (II, 1.3.2), donc
$H^q(f^{-1}(U),\mathcal F)=0$ par
(\hyperref[III.1.3.1-fr]{1.3.1}), ce qui démontre le corollaire.
\end{proof}

\begin{corollary}[1.3.3]
\label{III.1.3.3-fr}
Soient $Y$ un préschéma, $f:X\to Y$ un morphisme affine. Pour tout
$\mathcal O_X$-Module quasi-cohérent $\mathcal F$, l'homomorphisme canonique
$H^p(Y,f_*(\mathcal F))\to H^p(X,\mathcal F)$ (0, 12.1.3.1) est bijectif
pour tout $p$.
\end{corollary}

\begin{proof}
En effet, il suffit (en vertu de (0, 12.1.7)) de montrer que les
edge-homomorphismes
$''E_2^{p0}=H^p(Y,f_*(\mathcal F))\to H^p(X,\mathcal F)$ de la seconde suite
spectrale du foncteur composé $\Gamma f_*$ sont bijectifs. Mais le terme
$E_2$ de cette suite est donné par
$''E_2^{pq}=H^p(Y,R^qf_*(\mathcal F))$ (G, II, 4.17.1), donc il résulte de
(\hyperref[III.1.3.2-fr]{1.3.2}) que $''E_2^{pq}=0$ pour $q>0$, et la suite
spectrale dégénère~; d'où notre assertion (0, 11.1.6).
\end{proof}

\begin{corollary}[1.3.4]
\label{III.1.3.4-fr}
Soient $f:X\to Y$ un morphisme affine, $g:Y\to Z$ un morphisme.
